% =============================================================
%  APPENDIX D — K-FOLD EVALUATION PIPELINE AND RESULTS
% =============================================================

\section{K-Fold Evaluation Pipeline and Results}
\label{app:kfold}

% ------------------------------------------------------------------
\subsection{Experimental Protocol}
\label{app:kfold_protocol}
% ------------------------------------------------------------------

To obtain reliable and unbiased estimates of zero-shot classification accuracy, we adopt a stratified $K$-fold cross-validation protocol ($K=5$) applied uniformly to both the CLAP baseline and ResiDual$^*$\xspace. All experiments are conducted on four publicly available audio classification benchmarks: ESC-50~\cite{piczak2015esc}, IRMAS~\cite{bosch2012irmas}, TinySOL~\cite{cella2020tinysol}, and VocalSound~\cite{gong2022vocalsound}, which collectively span environmental sounds, musical instruments, orchestral timbres, and human vocal categories.

\paragraph{Dataset sampling.}
For each dataset, up to $N_{\mathrm{samples}} = 2{,}000$ examples are selected via class-stratified sampling with a fixed random seed, ensuring a balanced class distribution across all folds. The stratified sampling procedure first computes the per-class allocation $\lfloor N_{\mathrm{samples}} / C \rfloor$ (plus a unit remainder distributed to the first $N_{\mathrm{samples}} \bmod C$ classes), then samples without replacement from each class independently, guaranteeing reproducibility via a global seed.

\paragraph{Cross-validation split.}
The sampled subset is partitioned into $K=5$ stratified folds using \texttt{StratifiedKFold} with \texttt{shuffle=True} and a fixed \texttt{random\_state}. At each fold iteration, the training split (four folds) is further divided into two non-overlapping sub-partitions: a \emph{PCA fitting} partition of $N_{\mathrm{PCA}} = 250$ samples, used exclusively to estimate the spectral bases $\{\boldsymbol{\Phi}_{\ell,b,h},\, \boldsymbol{\mu}_{\ell,b,h}\}$ (cf.\ \S\ref{app:fitting}), and a \emph{$\lambda$-optimisation} partition comprising the remaining training examples, used to refine the weight vectors $\{\boldsymbol{\lambda}_{\ell,b,h}\}$ via the objective of Eq.~\eqref{eq:optim_obj}. The held-out fold (one fifth of the sampled subset) is reserved exclusively for final zero-shot evaluation.

\paragraph{Text prompt construction.}
Class-conditional text embeddings are computed once per dataset prior to
the fold loop. Each class label $c$ is formatted as the natural-language
prompt \texttt{"this is the sound of \{c\}"}, where underscores are
replaced by spaces. This prompt template follows the zero-shot inference
convention adopted in the official Microsoft CLAP
repository~\cite{microsoft_clap_github}, where the same phrasing is used
for audio classification evaluation across multiple benchmarks. The
resulting prompts are encoded by the frozen CLAP text encoder and
$\ell_2$-normalised to obtain $\mathbf{t}_c \in \mathbb{R}^{d_{\mathrm{proj}}}$
for each class $c \in \{1, \ldots, C\}$.
\paragraph{Zero-shot inference.}
At evaluation time, the audio embedding $\hat{\mathbf{a}}(x)$ for a test clip $x$ is obtained by passing the raw waveform through the (reweighted) audio encoder followed by the projection head $P$, and $\ell_2$-normalising the result. The predicted class is then
\begin{equation}
    \hat{y}(x) \;=\; \arg\max_{c \in \{1,\ldots,C\}}
    \hat{\mathbf{a}}(x)^\top \mathbf{t}_c.
    \label{eq:zs_inference}
\end{equation}
For the CLAP baseline, no adaptation is performed and the frozen model is evaluated directly. For ResiDual$^*$\xspace, the per-fold PCA fitting and $\lambda$-optimisation phases precede evaluation, and the model weights $\{\boldsymbol{\lambda}_{\ell,b,h}\}$ are fully reset between folds by reloading the CLAP checkpoint and re-instantiating all \texttt{SpectralReweightingLayer} modules.

\paragraph{Optimiser configuration.}
The weight vectors $\{\boldsymbol{\lambda}_{\ell,b,h}\}$ are updated with AdamW~\cite{loshchilov2019decoupled} with learning rate $\eta = 5 \times 10^{-3}$ and weight decay $\gamma = 5 \times 10^{-2}$, for a maximum of $T_{\max} = 40$ epochs with early stopping patience $p = 5$. The target stages are $\mathcal{L} = \{2, 3\}$ and the PCA variance threshold is $\rho_{\mathrm{var}} = 0.95$. All other hyperparameters follow Table~\ref{tab:residual_hparams}.

% ------------------------------------------------------------------
\subsection{Quantitative Results}
\label{app:kfold_results}
% ------------------------------------------------------------------

Table~\ref{tab:kfold_summary} reports the mean zero-shot accuracy and 
standard deviation across the five folds for each dataset and each model 
variant, comparing three ResiDual$^*$\xspace configurations that differ only in 
the set of target stages $\mathcal{L}$: reweighting applied to all stages 
($\mathcal{L}=\{0,1,2,3\}$), to the two deepest stages only 
($\mathcal{L}=\{2,3\}$, default), and to Stage~3 alone 
($\mathcal{L}=\{3\}$). 
Bold entries indicate the highest mean accuracy for each dataset across 
all variants. 
ResiDual$^*$\xspace with $\mathcal{L}=\{2,3\}$ improves over the CLAP baseline 
on all four benchmarks, with gains ranging from $+2.4$ percentage points 
on ESC-50 to $+54.8$ percentage points on TinySOL.

\begin{table*}[hb]
\centering
\caption{%
    Zero-shot accuracy ($\%$) under 5-fold cross-validation for CLAP 
    Classic and three ResiDual$^*$\xspace configurations differing in the set 
    of target stages $\mathcal{L}$.
    Mean $\pm$ standard deviation across folds.
    Bold indicates the highest accuracy for each dataset.
    $\Delta$ columns report the gain over CLAP Classic.
}
\label{tab:kfold_summary}
\setlength{\tabcolsep}{5pt}
\begin{tabular}{lcc ccc ccc}
\toprule
& & &
\multicolumn{2}{c}{$\mathcal{L}=\{3\}$} &
\multicolumn{2}{c}{$\mathcal{L}=\{2,3\}$ (default)} &
\multicolumn{2}{c}{$\mathcal{L}=\{0,1,2,3\}$} \\
\cmidrule(lr){4-5}\cmidrule(lr){6-7}\cmidrule(lr){8-9}
\textbf{Dataset} & \textbf{Classes} &
\textbf{CLAP Classic} &
\textbf{RD$^*$\xspace} & $\Delta$ &
\textbf{RD$^*$\xspace} & $\Delta$ &
\textbf{RD$^*$\xspace} & $\Delta$ \\
\midrule
ESC-50
    & 50
    & $93.85 \pm 1.06$
    & $95.45$ & $+1.60$
    & $\mathbf{96.20 \pm 0.43}$ & $+2.35$
    & $94.85$ & $+1.00$ \\
IRMAS
    & 11
    & $60.65 \pm 1.57$
    & $79.55$ & $+18.90$
    & $\mathbf{81.80 \pm 1.69}$ & $+21.15$
    & $79.70$ & $+19.05$ \\
TinySOL
    & 14
    & $35.55 \pm 0.64$
    & $73.34$ & $+37.79$
    & $\mathbf{90.31 \pm 1.76}$ & $+54.76$
    & $89.91$ & $+54.36$ \\
VocalSound
    & 6
    & $76.90 \pm 1.25$
    & $86.30$ & $+9.40$
    & $\mathbf{88.30 \pm 1.77}$ & $+11.40$
    & $87.10$ & $+10.20$ \\
\bottomrule
\multicolumn{9}{l}{%
    \textsuperscript{$\dagger$}TinySOL contains 1\,913 valid samples 
    after validation; all are used.
    Std.\ deviations omitted for $\mathcal{L}\neq\{2,3\}$ for brevity.}
\end{tabular}
\end{table*}

\paragraph{Effect of target stage selection.}
The comparison across the three configurations in 
Table~\ref{tab:kfold_summary} reveals a clear and consistent pattern: 
restricting reweighting to the two deepest stages ($\mathcal{L}=\{2,3\}$) 
yields the best results on all four benchmarks, while extending it to all 
stages ($\mathcal{L}=\{0,1,2,3\}$) or restricting it to Stage~3 alone 
($\mathcal{L}=\{3\}$) both lead to lower accuracy.

The underperformance of $\mathcal{L}=\{0,1,2,3\}$ relative to the default 
$\mathcal{L}=\{2,3\}$ is particularly informative. 
As established by the head analysis in Appendix~\ref{app:head_analysis} 
(Figures~\ref{fig:blockwise_metrics}--\ref{fig:head_metrics_panels}), 
early-stage heads (Stages~0 and~1) operate in near-unidimensional regimes 
with very high PC1 dominance ($\mathrm{EVR}_1 \approx 0.92$ in Stage~0) 
and near-zero semantic coherence Z-scores, indicating that they encode 
low-level acoustic properties --- such as broadband energy and dominant 
spectral envelope --- that are common across sound categories rather than 
discriminative for them. 
Applying spectral reweighting to these heads therefore introduces 
unnecessary degrees of freedom into the optimisation, risking interference 
with generic low-level feature extraction that the downstream stages rely 
upon, without providing additional discriminative signal. 
The result is a modest but consistent degradation relative to reweighting 
only Stages~2 and~3, where semantic specialisation has already emerged.

The configuration $\mathcal{L}=\{3\}$ provides a useful ablation showing 
that Stage~2 contributes meaningfully to the adaptation gain. 
Stage~2 is precisely the regime identified in Appendix~\ref{app:head_analysis} 
as the onset of functional diversification, where intrinsic nonlinear 
dimensionality peaks ($L/N$ ratio at its maximum of $2.65$) and the first 
semantically coherent heads appear 
(e.g.\ \texttt{L2\_B4\_H1}, $Z=1.46$; \texttt{L2\_B5\_H15}, $Z=1.23$). 
Excluding Stage~2 from the reweighting thus forfeits the ability to amplify 
these early-emerging discriminative directions, with the loss being 
especially pronounced on TinySOL ($-16.97$\,pp relative to 
$\mathcal{L}=\{2,3\}$), the benchmark most reliant on fine-grained timbral 
discrimination.

Taken together, these results validate the stage-selection criterion 
proposed in \S\ref{sec:method}: targeting the stages where semantic 
coherence is highest and nonlinear manifold complexity is maximal, as 
identified by the head analysis, is both necessary and sufficient for 
optimal performance.

Table~\ref{tab:kfold_folds} provides the per-fold breakdown for the 
default configuration $\mathcal{L}=\{2,3\}$, offering a more granular 
view of stability across splits.

\begin{table*}[ht]
\centering
\caption{%
    Per-fold zero-shot accuracy ($\%$) for CLAP Classic and 
    ResiDual$^*$\xspace ($\mathcal{L}=\{2,3\}$).
    Bold entries indicate the higher accuracy within each fold and dataset.
}
\label{tab:kfold_folds}
\newcolumntype{L}{>{\raggedright\arraybackslash}X}
\newcolumntype{C}{>{\centering\arraybackslash}X}
\setlength{\tabcolsep}{5pt}
\begin{tabularx}{\textwidth}{LLCCCCC}
\toprule
\textbf{Dataset} & \textbf{Model} &
\textbf{Fold 1} & \textbf{Fold 2} & \textbf{Fold 3} &
\textbf{Fold 4} & \textbf{Fold 5} \\
\midrule
\multirow{2}{*}{ESC-50}
    & CLAP Classic
        & $92.25$ & $95.00$ & $93.75$ & $93.25$ & $95.00$ \\
    & ResiDual$^*$\xspace
        & $\mathbf{95.50}$ & $\mathbf{96.75}$
        & $\mathbf{96.50}$ & $\mathbf{96.00}$ & $\mathbf{96.25}$ \\
\midrule
\multirow{2}{*}{IRMAS}
    & CLAP Classic
        & $59.50$ & $60.00$ & $62.75$ & $62.25$ & $58.75$ \\
    & ResiDual$^*$\xspace
        & $\mathbf{84.50}$ & $\mathbf{80.50}$
        & $\mathbf{83.00}$ & $\mathbf{81.00}$ & $\mathbf{80.00}$ \\
\midrule
\multirow{2}{*}{TinySOL}
    & CLAP Classic
        & $35.24$ & $35.24$ & $34.96$ & $35.53$ & $36.78$ \\
    & ResiDual$^*$\xspace
        & $\mathbf{87.68}$ & $\mathbf{92.26}$
        & $\mathbf{92.26}$ & $\mathbf{89.97}$ & $\mathbf{89.37}$ \\
\midrule
\multirow{2}{*}{VocalSound}
    & CLAP Classic
        & $74.75$ & $77.00$ & $76.50$ & $78.00$ & $78.25$ \\
    & ResiDual$^*$\xspace
        & $\mathbf{88.00}$ & $\mathbf{90.75}$
        & $\mathbf{85.50}$ & $\mathbf{87.75}$ & $\mathbf{89.50}$ \\
\bottomrule
\end{tabularx}
\end{table*}

% ------------------------------------------------------------------
\subsection{Summary}
\label{app:kfold_summary_sec}
% ------------------------------------------------------------------

The optimisation dynamics are consistent with the viability of the RD$^*$\xspace
pipeline: convergence is broadly monotone in the early training phase, and
early stopping reliably identifies configurations that substantially outperform
the frozen baseline on held-out data.
The gap between optimisation-partition and test accuracy is small on ESC-50
but more pronounced on IRMAS and VocalSound, pointing to a generalisation
challenge that scales with task difficulty rather than to overfitting of the
$\lambda$ parameters per se.
These results collectively validate the design choices described in
Appendix~\ref{app:residual_impl}: a fixed PCA basis estimated on a small
unlabelled subset, combined with a lightweight per-component rescaling
optimised via cross-entropy on a modest labelled partition, suffices to
substantially improve zero-shot classification accuracy without modifying
any pretrained parameters, even under conditions of imperfect
optimisation-to-test generalisation.